\section{主线确认：收盘后只查五件事}

把D+1和D+2当成考试日。五项都过，题材才从“看看”升级为“继续跟踪”；少一项，结论回到观察或结束。核心原则很简单：主线需要持续的消息、可交易的梯队和主板承接，单日大涨不够。

\begin{exhibitbox}[展品：五项检查表]
\begin{tabularx}{\textwidth}{L{2.2cm}X X}
\toprule
检查项 & 通过的样子 & 失败时怎么做 \\
\midrule
催化是否续上 & 出现官方平台、制作方、监管或公司资料；内容包含播出反馈或第二个可追溯事件。 & 继续观察，题材不升级。 \\
平台锚点是否还在 & 芒果TV页面敞口仍可核验；替代锚点须有公司级关系资料。 & 删除“换龙头”想象。 \\
主板梯队是否出现 & 至少一只项目直接关系主板股走出1到2板，同时有同层首板补位。 & 主线假设结束。 \\
主板容量是否承接 & 项目直接关系主板股连续两日成交额靠前，收盘不弱于前收，全天未被一字锁住。 & 只看广义AI内容扩散。 \\
市场是否配合 & 合格集合至少三只，至少三分之二收涨；市场数据未给出风险释放或执行阻断。 & 停止升级计划。 \\
\bottomrule
\end{tabularx}
\end{exhibitbox}

当前第三、第四项无法通过，因为主板的项目直接关系公司尚未出现。这也是报告最重要的结论：题材可以热，龙头战法仍然要等结构。
